% !TEX TS-program = lualatex
\documentclass{book}
\usepackage{amsmath}
\usepackage{graphicx}
\input{shorthands}

\newcommand{\vect}[1]{\overrightarrow{#1}}
\newcommand{\unit}[2]{#1\,\mathrm{#2}}
\newenvironment{remark}{\par\medskip\noindent\textbf{Remark.}\ }{\par\medskip}

\begin{document}

\chapter{Point mechanics}

\section{Positions and velocities}

A point is given by $\vect{r}(t)$, its velocity by $\vect{v} = \dot{\vect{r}}$.
The speed of a free fall after $t = \unit{2}{s}$ is about $\unit{20}{m.s^{-1}}$.

\begin{itemize}
\item the position, $\vect{r}$;
\item the velocity, $\vect{v}$;
\item the acceleration, $\vect{a}$.
\end{itemize}

\begin{figure}
\centering
\includegraphics[width=0.6\textwidth]{figures/trajectory}
\caption{A trajectory in the plane.}
\end{figure}

\section{Energy}

The kinetic energy of a point of mass $m$ is
\[
    E_c = \frac{1}{2} m v^2 .
\]

\begin{center}
Energy is conserved when no non-conservative force does any work.
\end{center}

\begin{remark}
The radius runs over $r\in\[0;\infty\[$, written with escaped brackets: inside
$\dots$ those are commands, not the display math of \verb|\[ \]|.
\end{remark}

\begin{figure}
\centering
\includegraphics[width=0.5\textwidth]{figures/energy}
\caption{Potential energy against the radius.}
\end{figure}

\section{Exercises}

\begin{enumerate}
\item Show that $\vect{a}$ is normal to $\vect{v}$ for a uniform circular motion.
\item Give $E_c$ for $v = \unit{3}{m.s^{-1}}$ and $m = \unit{2}{kg}$.
\end{enumerate}

\begin{center}
\includegraphics[width=0.3\textwidth]{figures/circle}
\end{center}

\end{document}
